\chapter{Ambiente de simulação e redes de teste}
\label{cap:ambiente}

Este capítulo descreve a infraestrutura computacional desenvolvida para a execução das simulações e as 
redes de distribuição utilizadas. A rede IEEE~33 barras foi empregada para a validação inicial do 
ambiente, ao passo que a rede IEEE~34 barras, de maior porte e com elementos de regulação 
representativos, foi adotada como sistema de estudo nas etapas subsequentes.

\section{O software OpenDSS}
\label{sec:opendss}

O OpenDSS (\emph{Open Distribution System Simulator}), desenvolvido pelo 
\emph{Electric Power Research Institute} (EPRI)~\cite{dugan2020opendss}, é um simulador de propósito 
geral para análise de sistemas de distribuição, amplamente utilizado para o cálculo de fluxo de 
potência. A matriz de admitância do sistema, que relaciona correntes injetadas e tensões nodais, é 
montada a partir das matrizes de admitância primitivas de cada elemento e resolvida por um \emph{solver} 
de matrizes esparsas. O fluxo de potência, que determina as tensões em cada barra e os fluxos de 
potência ativa e reativa nas linhas, é calculado iterativamente pelo método de injeção de corrente 
(padrão), com o método de Newton disponível para circuitos de convergência mais difícil. Determinar o 
fluxo de potência é fundamental para avaliar o carregamento da rede, identificar violações de tensão e 
apoiar as decisões de controle.

\section{Integração Python--OpenDSS}
\label{sec:api}

A integração do OpenDSS com Python possibilita o desenvolvimento de ferramentas que exploram os 
resultados do fluxo de potência de forma programática e automatizada, viabilizando estudos de 
controladores de BESS em diferentes cenários. Para essa integração, avaliaram-se as interfaces de 
programação disponíveis. A abordagem oficial baseia-se na interface \emph{Component Object Model} 
(COM), de uso restrito ao Windows. Como alternativa multiplataforma, o projeto DSS-Extensions 
disponibiliza três bibliotecas Python --- DSS-Python, OpenDSSDirect.py e AltDSS-Python ---, todas 
construídas sobre o motor AltDSS, compatível com o OpenDSS oficial e com suporte a Windows, Linux e 
macOS.

Cabe registrar um ajuste em relação ao plano de trabalho, que previa a interface COM: optou-se pela 
biblioteca OpenDSSDirect.py. A escolha justifica-se por ser o pacote mais maduro entre as opções do 
DSS-Extensions, por oferecer uma interface por chamadas de função direta e bem documentada, adequada 
às iterações frequentes sobre elementos da rede exigidas pelas etapas de controle, e por disponibilizar 
funções utilitárias para conversão dos elementos da rede em estruturas de dados da biblioteca 
\emph{pandas}, facilitando a exportação e a análise dos resultados. Além disso, o suporte 
multiplataforma elimina a dependência do sistema operacional Windows.

\section{Validação na rede IEEE~33 barras}
\label{sec:ieee33}

A rede IEEE~33 barras é um sistema de distribuição radial de referência, proposto originalmente por 
Baran e Wu~\cite{baranwu1989}, com tensão nominal de 12,66~kV (média tensão), composto por 33 barras e 
32 ramais. A topologia é puramente radial, com uma barra de suprimento (barra~1, tomada como referência) 
e 32 barras de carga; a demanda total é de 3,715~MW de potência ativa, distribuída de forma não 
uniforme. As \emph{tie lines} originais, com chaves normalmente abertas, foram desconsideradas nesta 
simulação. A Figura~\ref{fig:ieee33} apresenta o diagrama unifilar da rede.

\begin{figure}[!htb]
  \centering
  \caption{Diagrama unifilar do sistema de testes IEEE~33 barras.}
  \label{fig:ieee33}
  \includegraphics[width=0.92\textwidth]{diagrama_unifilar_33_barras.jpg}
  \par\vspace{2pt}\fonte{adaptado de~\cite{baranwu1989}}
\end{figure}

Para validar o ambiente, executou-se o fluxo de potência da rede em modo estático (\emph{snapshot}), 
com a rede operando a 100\% de sua carga nominal. Os resultados revelaram um perfil de tensão 
decrescente ao longo dos alimentadores, comportamento esperado em redes radiais com carga distribuída: 
a barra de suprimento apresenta tensão próxima da nominal, enquanto as barras mais distantes da 
subestação atingem valores inferiores ao limite de 0,95~p.u., com mínimo de 0,9059~p.u. na barra~18. 
As tabelas completas de tensões nodais e de fluxos de potência por linha encontram-se no 
Apêndice~\ref{ap:fluxo33}. A correta obtenção desses resultados validou o ambiente de cossimulação e 
estabeleceu a infraestrutura necessária para as etapas de controle. Cabe registrar que a rede IEEE~33 
barras foi empregada apenas nessa validação inicial; melhorias propostas na literatura, como a inclusão 
de geração distribuída, compensadores de reativos e curvas de carga~\cite{ghorbanian2019}, motivaram a 
adoção da rede IEEE~34 barras como sistema de estudo nas etapas seguintes.

\section{Rede de estudo: IEEE~34 barras}
\label{sec:ieee34}

A partir da validação do ambiente, adotou-se a rede \emph{IEEE 34 Bus Test Feeder} como sistema de 
estudo, por ser um alimentador de média tensão amplamente utilizado na avaliação de algoritmos de 
controle, fluxo de potência e integração de recursos distribuídos, e por incorporar dispositivos de 
regulação representativos de redes reais. A Figura~\ref{fig:ieee34} apresenta o diagrama da rede.

\begin{figure}[!htb]
  \centering
  \caption{Diagrama da rede IEEE~34 barras, com indicação das categorias de carga.}
  \label{fig:ieee34}
  \includegraphics[width=\textwidth]{diagrama_unifilar_34_barras.png}
  \par\vspace{2pt}\fonte{autoria própria}
\end{figure}

A alimentação é feita a partir de uma subestação conectada à barra principal por um transformador 
abaixador trifásico de 25~MVA e relação 69/24,9~kV, com impedância de curto-circuito reduzida, de modo 
a tornar a fonte suficientemente rígida para os estudos de penetração de GD. A rede opera 
predominantemente em 24,9~kV, com exceção de um ramal (barras 888 e 890) alimentado por um segundo 
transformador de 500~kVA e relação 24,9/4,16~kV. A topologia é radial, com trechos trifásicos e 
monofásicos; as 32 linhas são definidas a partir dos códigos de impedância de fase padronizados pelo 
IEEE.

Para o controle de tensão ao longo do alimentador, a rede possui dois bancos reguladores com controle 
independente por fase: o primeiro, entre as barras 814 e 814R, opera com tensão de referência de 122~V 
e banda de 2~V; o segundo, entre as barras 852 e 852R, é ajustado para 124~V. Ambos foram modelados 
como transformadores monofásicos de enrolamento duplo com controle automático de tap via o elemento 
\emph{RegControl} do OpenDSS. Há ainda dois bancos de capacitores \emph{shunt} nas barras 844 e 848, 
que contribuem para o suporte de tensão nas regiões mais afastadas da subestação. As cargas dividem-se 
em concentradas (\emph{spot loads}), nas barras 860, 840, 844, 848, 830 e 890, e distribuídas, 
alocadas nas barras de origem e destino de cada segmento; a carga de pico total é de 1,769~MW e a 
demanda reativa total, de 1,044~MVAr.

\subsection{Perfis de carga}
\label{sec:perfis_carga}

À rede original foram acrescentados três perfis de carga horária --- residencial, comercial e industrial 
--- com resolução de uma hora e horizonte de 24~horas, representando o comportamento típico de cada 
categoria de consumidor. Cada carga foi então associada a um desses perfis, de modo que a demanda a 
cada hora é obtida multiplicando o fator de carga horário pela demanda máxima da carga. 
A Figura~\ref{fig:perfis_carga} ilustra os três perfis. O perfil residencial caracteriza-se por 
consumo reduzido na madrugada e pico por volta das 20~h; o comercial apresenta consumo elevado entre 
9~h e 12~h, cedendo à noite; e o industrial distingue-se dos demais por apresentar fatores
multiplicativos superiores à unidade entre 9~h e 12~h, indicando que a potência nominal declarada
corresponde a uma condição de carga parcial, superada durante o turno de operação plena. Foram 
classificadas como residenciais as cargas das barras 848 e 860; como comerciais, as das barras 820, 
822, 824, 826, 832, 836, 840 e 858; e como industriais, as das barras 830, 834, 844 e 890. As demais 
cargas distribuídas, de menor expressão, foram tratadas como residenciais.

\begin{figure}[!htb]
  \centering
  \caption{Curvas de carga típicas (residencial, comercial e industrial).}
  \label{fig:perfis_carga}
  \includegraphics[width=0.75\textwidth]{curvas_carga.png}
  \par\vspace{2pt}\fonte{autoria própria}
\end{figure}
